\documentclass[conference]{IEEEtran}

\usepackage{amsmath,amssymb}
\usepackage{graphicx}
\usepackage{booktabs}
\usepackage{hyperref}
\usepackage{xcolor}
\usepackage{microtype}

\hypersetup{
  colorlinks=true,
  linkcolor=blue,
  citecolor=blue,
  urlcolor=blue
}

% ─────────────────────────────────────────────────────────────────────────────
\title{Estimating Retail Promotion Spillovers\\with Graph-Aware Double Machine Learning}

\author{
  \IEEEauthorblockN{David M. McKinney}
  \IEEEauthorblockA{Independent Researcher\\
    Ann Arbor, MI, USA}
}

\begin{document}
\maketitle

% ─────────────────────────────────────────────────────────────────────────────
\begin{abstract}
    Promotional pricing decisions in retail don't just affect the promoted
    product --- they ripple across substitute and complementary items in ways
    that standard causal estimators miss entirely.
    We propose a \emph{graph-aware Double Machine Learning} framework that
    explicitly models these spillovers by building a product co-purchase
    and similarity graph, computing weighted neighborhood exposure variables,
    and fitting a linear DML estimator with two-way cluster-robust standard
    errors.
    On the Dunnhumby retail panel (49.9\,M transactions, 92\,K products),
    we estimate a direct promotion effect of $+0.037$ on the
    $\log(1 + \text{units sold})$ scale
    (95\% CI: $[0.034,\, 0.041]$), a statistically significant substitute
    cannibalization effect of $-0.0022$ ($p < 0.01$), and a positive but
    noisy complement attachment effect.
    An ablation study shows that removing Node2Vec product embeddings inflates
    the direct ATE by 16\% and the substitute spillover by 2.2$\times$,
    with a shift consistent with embeddings capturing latent demand structure
    and product-level omitted variables that simpler approaches leave behind.
    A semi-synthetic evaluation finds low bias overall, but also shows that
    direct-effect confidence intervals undercover, so the uncertainty on the
    direct effect should be interpreted cautiously.
\end{abstract}

\begin{IEEEkeywords}
    causal inference, double machine learning, graph spillovers,
    retail pricing, treatment effects, exposure mapping
\end{IEEEkeywords}


% ─────────────────────────────────────────────────────────────────────────────
\section{Introduction}

When a grocery retailer puts yogurt on promotion, sales go up --- but
which yogurts? Customers who were going to buy a competing brand may switch
over (substitution), while shoppers buying yogurt may also grab granola
(complementarity).
Standard causal estimators treat each product in isolation and miss both
effects.

Capturing these spillovers is not just theoretically interesting; it has
real business consequences.
A pricing team that ignores cannibalization will over-count the value of
promotions, and one that ignores complement attachment will under-invest in
category-level cross-promotion strategies.

This paper makes three main contributions:

\begin{enumerate}
    \item We formalize retail pricing spillovers through the lens of
          \emph{exposure mapping} \cite{aronow2017estimating}, representing each
          product's neighborhood treatment intensity as a weighted aggregate of
          its graph neighbors' promotion statuses.

    \item We integrate these exposure variables into a
          \emph{Double ML} \cite{chernozhukov2018double} framework with
          LightGBM \cite{ke2017lightgbm} nuisance models and two-way clustered
          standard errors \cite{cameron2011robust}, yielding scalable
          cluster-robust inference over
          49.9\,M retail transactions.

    \item Through a five-way ablation study, we demonstrate that
          Node2Vec \cite{grover2016node2vec} product embeddings are the most
          critical modeling choice: removing them biases estimates substantially,
          while the choice of binary vs.\ edge-weighted adjacency has almost no
          impact.
\end{enumerate}

Our central empirical point is not just that graphs add predictive features.
By encoding which products interact, graph structure lets us estimate
downstream treatment effects that standard no-graph specifications either
omit or misattribute.


% ─────────────────────────────────────────────────────────────────────────────
\section{Related Work}

\paragraph{Double Machine Learning.}
Chernozhukov et al.\ \cite{chernozhukov2018double} introduced DML to
remove regularization bias when using flexible machine learning models for
nuisance estimation.
Their cross-fitting procedure guarantees $\sqrt{n}$-consistent, Gaussian
inference under weak conditions.
We adopt their \texttt{LinearDML} implementation from
EconML \cite{battocchi2019econml}.

\paragraph{Interference and Spillovers.}
The SUTVA assumption of standard causal inference---that a unit's outcome
depends only on its own treatment---fails whenever units interact.
Hudgens and Halloran \cite{hudgens2008toward} laid the groundwork for
causal inference under interference, and Aronow and
Samii \cite{aronow2017estimating} introduced exposure mapping as a general
device for specifying how spillovers enter the potential outcomes model.
We apply exposure mapping to the product-level retail setting, where
``interference'' is mediated by the product graph.

\paragraph{Heterogeneous Treatment Effects.}
Wager and Athey \cite{wager2018estimation} developed causal forests for
non-parametric CATE estimation.
We use a causal forest as one baseline but rely on the linear DML as our
primary estimator, since interpretability and valid standard errors are more
important than flexibility for our spillover decomposition.

\paragraph{Product Graphs in Retail.}
Representing product relationships as a graph and using graph neural
networks for recommendation has been studied extensively.
Our work is different: we use the graph not for prediction but to construct
valid exposure variables for causal estimation, keeping graph structure
inside the identification argument rather than treating it as a
feature-engineering step.


% ─────────────────────────────────────────────────────────────────────────────
\section{Data}

\subsection{The Dunnhumby Complete Journey Dataset}

We use the publicly available Dunnhumby ``Complete Journey''
dataset \cite{dunnhumby2020journey}, which records household-level
purchases at a US grocery chain over roughly two years.
The raw dataset contains 49.9\,M transaction records spanning 2,500
households, 391 stores, and 92,339 unique products across 44 departments.

We aggregate to a \emph{panel} at the grain
$(i, s, t) = (\text{PRODUCT\_ID}, \text{STORE\_ID}, \text{WEEK\_NO})$.
The primary outcome is \texttt{units\_sold}, log-transformed as
$Y = \log(1 + \text{units})$ to handle the zero-inflated count
distribution.
Products with mean weekly sales below 0.01 units are dropped as structural
zeros, leaving 31,142 active products and 49.9\,M panel rows.
Experiments use a 6\,M-row random subsample of the training split for
computational tractability.

\subsection{Substitute Graph}

We train Node2Vec \cite{grover2016node2vec} embeddings on the
co-purchase graph to obtain 64-dimensional product representations.
The substitute adjacency matrix $A^{\mathrm{sub}}$ is built by computing
cosine similarity between embedding pairs and retaining the top-$k$
neighbors above a similarity threshold.
Edge weights reflect semantic product similarity.

\subsection{Complement Graph}

The complement adjacency matrix $A^{\mathrm{comp}}$ is derived from
Positive Pointwise Mutual Information (PPMI) \cite{church1990word}
applied to basket-level co-purchase counts.
High PPMI between two products indicates they are purchased together more
often than chance, capturing complementarity.

\subsection{Treatment and Covariates}

The primary treatment variable is \texttt{promo\_any}, a binary indicator
for whether a product received any display or mailer promotion in a given
store-week.
A secondary model estimates effects of the continuous treatment
\texttt{discount\_rate} (fractional price markdown).
The covariate matrix $W$ includes store ID, week number, department,
commodity description, sub-commodity, brand, and Node2Vec embedding
coordinates for each product.
These controls are intended to absorb the main observable drivers of both
promotion assignment and demand: store ID captures persistent
location-level merchandising differences, week number captures calendar and
seasonal demand shifts, product taxonomy captures broad assortment
structure, and embedding coordinates provide dense product-level controls
for latent similarity in demand and promotion propensity.

\subsection{Experimental Setup}

All reported graph-DML estimates use the temporal training window from
weeks 1--70, with weeks 71--85 reserved for validation and weeks 86--102
reserved for test-time extensions.
The published graph-DML, discount, and ablation tables are fit on a
6\,M-row random subsample of the training split, which drops to
5.19\,M rows after filtering products with mean weekly demand below 0.01.

The complement graph retains product pairs with at least 10 same-basket
co-occurrences and weights retained edges by normalized PPMI.
Node2Vec embeddings use dimension 64, 10 walks per node, walk length 80,
$p = q = 1$ (the unbiased DeepWalk setting), 5 Word2Vec epochs, and seed 42.
The substitute graph connects products within the same
\texttt{SUB\_COMMODITY\_DESC} and keeps up to 50 neighbors per node,
reweighted by cosine similarity in embedding space.

For nuisance estimation, both $E[Y \mid W]$ and $E[T \mid W]$ are modeled
with LightGBM regressors using 200 trees, depth 6, learning rate 0.1,
31 leaves, minimum child sample size 20, and 5 threads.
Cross-fitting uses $K=5$ folds via \texttt{GroupKFold}, grouping on
\texttt{PRODUCT\_ID} so that all store-weeks for a product stay in the same
fold.
The primary outcome is modeled on the $\log(1 + \text{units sold})$ scale.
For the continuous-treatment specification, \texttt{discount\_rate} is scaled
to $[0,1]$, so its coefficients should be interpreted per unit change in the
fractional discount rate rather than per one-percentage-point change.


% ─────────────────────────────────────────────────────────────────────────────
\section{Methodology}

\subsection{Exposure Mapping}

For each panel observation $(i, s, t)$, we define two spillover treatment
variables: the weighted fraction of substitute and complement neighbors that
are currently promoted.

\begin{equation}
    e_i^{\mathrm{sub}} = \frac{\sum_{j} A^{\mathrm{sub}}_{ij}\, T_{jst}}%
    {\sum_{j} A^{\mathrm{sub}}_{ij}\, \mathbf{1}[j \in \mathcal{S}_{st}]},
    \label{eq:exposure}
\end{equation}

\noindent and analogously for $e_i^{\mathrm{comp}}$ using $A^{\mathrm{comp}}$.
Here $\mathcal{S}_{st}$ is the set of products available in store $s$ at
week $t$, and $T_{jst}$ is the promotion indicator for product~$j$.
The denominator normalizes by the number of observable neighbors, not the
total degree, to avoid boundary effects from sparse store assortments.

This one-hop exposure mapping assumes that spillovers propagate only through
direct graph neighbors.
That assumption is substantively reasonable in grocery retail, where the
largest substitution and attachment responses are typically concentrated in
the closest product alternatives and co-purchase partners.
Accordingly, our estimates should be interpreted as \emph{local one-hop}
spillover effects induced by the observed product graph.
Expanding the neighborhood radius is a natural robustness check, but we do
not pursue it here.

\subsection{Double ML Estimator}

We model the bivariate treatment as
$\mathbf{T}_i = [T_i,\, e_i^{\mathrm{sub}},\, e_i^{\mathrm{comp}}]^\top$
and estimate:

\begin{equation}
    Y_i = \boldsymbol{\theta}^\top \mathbf{T}_i + g(W_i) + \epsilon_i,
    \label{eq:model}
\end{equation}

\noindent where $g(\cdot)$ is a nuisance function approximated by
LightGBM \cite{ke2017lightgbm} and $\boldsymbol{\theta}$ contains
the three causal effects of interest.

We use the \texttt{LinearDML} implementation from EconML
\cite{battocchi2019econml} with $K=5$ cross-fitting folds, where fold
assignment groups all store-weeks of a given product together
(\texttt{GroupKFold}) to prevent leakage between nuisance and final-stage
estimation.

\subsection{Identification Assumptions}

Our causal interpretation relies on a conditional ignorability assumption for
both own treatment and graph exposure: conditional on $W_i$, the joint
treatment vector $[T_i, e_i^{\mathrm{sub}}, e_i^{\mathrm{comp}}]$ is assumed
as-good-as-random with respect to potential outcomes.
Operationally, this means that after conditioning on product identity and
similarity proxies (taxonomy and embeddings), store effects, and calendar
time, there are no remaining systematic drivers of promotion assignment that
also directly shift demand.

Under this design, $W$ is not just a predictive feature set; it is the set of
observed controls used to block the main backdoor paths from promotional
assignment to sales.
The DML procedure then flexibly partials out $E[Y \mid W]$ and
$E[T \mid W]$ before estimating the low-dimensional treatment effects.
DML is critical here because spillover-adjusted treatment effects are
estimated in a high-dimensional setting where naive regularization of the
nuisance models could otherwise bias the target coefficients.
We also require overlap: products with comparable observed characteristics
must have non-degenerate variation in promotion and neighborhood exposure.

Because the data are observational, this assumption is ultimately
unverifiable.
Residual confounding may remain if unobserved store-week shocks---such as
local events, competitor actions, or supplier-funded campaigns---jointly
affect promotion decisions and demand for clusters of related products.
Our results should therefore be read as estimates under conditional
ignorability given the observed product, store, time, and embedding-based
controls.

\subsection{Two-Way Cluster-Robust Standard Errors}

Errors are correlated within products (same item promoted across stores)
and within stores (chain-wide promotional events).
We apply the Cameron-Gelbach-Miller (CGM) two-way clustering
estimator \cite{cameron2011robust}:

\begin{equation}
    \hat{V}_{\mathrm{2way}} = \hat{V}_{\mathrm{product}}
    + \hat{V}_{\mathrm{store}}
    - \hat{V}_{\mathrm{product}\times\mathrm{store}},
    \label{eq:cgm}
\end{equation}

\noindent where each one-way variance is computed using the Liang-Zeger
sandwich estimator \cite{liang1986longitudinal} applied to the
DML residuals.
The final-stage OLS uses a pseudoinverse ($\hat{X}^+$) rather than a
standard inverse to gracefully handle ablation conditions where one
treatment column is identically zero.


% ─────────────────────────────────────────────────────────────────────────────
\section{Results}

\subsection{Main Effects}

Table~\ref{tab:main} shows the estimated average treatment effects for the
primary binary-promotion model.

\begin{table}[ht]
    \centering
    \caption{Main causal effect estimates (binary promotion treatment).
        95\% confidence intervals use two-way cluster-robust SEs clustered
        on product and store. $n = 5.19\,\mathrm{M}$ observations.}
    \label{tab:main}
    \begin{tabular}{lccc}
        \toprule
        Effect               & ATE       & CI lower  & CI upper  \\
        \midrule
        Direct (own promo)   & $+0.0373$ & $0.0337$  & $0.0408$  \\
        Substitute spillover & $-0.0022$ & $-0.0036$ & $-0.0007$ \\
        Complement spillover & $+0.0003$ & $-0.0010$ & $+0.0016$ \\
        \bottomrule
    \end{tabular}
\end{table}

The direct promotion effect is positive and precisely estimated on the
$\log(1 + \text{units sold})$ scale: the direct coefficient is 0.037 with a
tight clustered confidence interval, corresponding to roughly a 3.8\%
increase in expected units as a back-of-the-envelope translation.
The substitute spillover is negative and statistically significant at
$p < 0.01$, confirming cannibalization: when a product's graph substitutes
are promoted, the focal product's expected outcome declines modestly.
The complement spillover is positive but small and does not reach
conventional significance levels, suggesting that complement attachment lift
exists directionally but is too noisy to detect reliably in this dataset.

\subsection{Baseline Comparison}

Table~\ref{tab:baselines} benchmarks our graph-aware DML against four
alternative estimators on the direct ATE.
Figure~\ref{fig:baselines} plots these comparisons visually.

\begin{table}[ht]
    \centering
    \caption{Direct ATE estimates across estimators.
        OLS and DML no-graph use no spillover features;
        DML unweighted spillover uses binary co-purchase fractions
        rather than the graph-weighted exposure variables.}
    \label{tab:baselines}
    \begin{tabular}{lcc}
        \toprule
        Model                       & ATE                & 95\% CI              \\
        \midrule
        OLS (no graph)              & $+0.0404$          & $[0.0401,\, 0.0406]$ \\
        Causal forest (no graph)    & $+0.0307$          & $[0.0152,\, 0.0462]$ \\
        DML (no graph)              & $+0.0290$          & $[0.0285,\, 0.0295]$ \\
        DML + unweighted spillover  & $+0.0270$          & $[0.0252,\, 0.0289]$ \\
        \textbf{DML + graph (ours)} & $\mathbf{+0.0373}$ & $[0.0337,\, 0.0408]$ \\
        \bottomrule
    \end{tabular}
\end{table}

OLS without graph features produces the largest direct ATE ($+0.040$) ---
its tight confidence interval is an artifact of ignoring non-linear
confounding.
The standard DML estimate ($+0.029$) is lower once flexible confound
adjustment is applied.
Adding unweighted spillover controls reduces the direct effect further
($+0.027$), consistent with spillover confounding in the no-graph estimates.
Our graph-aware DML yields a slightly higher direct effect ($+0.037$)
than the no-graph DML baselines.
This pattern is consistent with spillover structure affecting which observed
variation is attributed to own-treatment versus neighbor-treatment channels,
although the real-data setting does not identify the ground-truth
decomposition.

\begin{figure}[t]
    \centering
    \includegraphics[width=\linewidth]{../results/figures/baselines_vs_graph_direct.png}
    \caption{Direct ATE with 95\% CIs across five estimators. Error bars
        for the causal forest are wider due to the 10\,M-row subsample used
        for that estimator.}
    \label{fig:baselines}
\end{figure}

\subsection{Discount Rate Model}

Table~\ref{tab:discount} shows results for the continuous
\texttt{discount\_rate} treatment.

\begin{table}[ht]
    \centering
    \caption{Estimated effects in the continuous discount-rate model.
        The treatment is \texttt{discount\_rate} on the $[0,1]$ scale, so
        coefficients are reported per unit change in fractional discount depth.}
    \label{tab:discount}
    \begin{tabular}{lccc}
        \toprule
        Effect               & ATE      & CI lower & CI upper \\
        \midrule
        Direct               & $+2.43$  & $+2.38$  & $+2.47$  \\
        Substitute spillover & $-0.016$ & $-0.046$ & $+0.015$ \\
        Complement spillover & $+0.014$ & $+0.006$ & $+0.021$ \\
        \bottomrule
    \end{tabular}
\end{table}

Coefficients are reported on the $\log(1 + \text{units sold})$ scale per
unit change in \texttt{discount\_rate}, where a unit corresponds to moving
from a 0\% to 100\% discount.
Under a local linear approximation, a 10-percentage-point deeper discount
corresponds to about 0.243 log-points on the direct channel.
The direct effect is strongly positive, indicating meaningful price
sensitivity beyond the binary promotion indicator.
The complement spillover is now statistically significant ($p < 0.01$):
deeper discounting on a product generates measurable attachment lift for its
graph complements, consistent with customers expanding their basket when a
core item is more deeply discounted.
Taken together with the near-zero complement estimate in the binary-promotion
model, this suggests that coarse promotion indicators may be too weak to
detect complementarity, whereas continuous discount depth carries a stronger
behavioral signal.
For substitutes, the spillover is negative but does not reach significance
under continuous treatment.


% ─────────────────────────────────────────────────────────────────────────────
\section{Heterogeneous Treatment Effects}

\begin{figure}[t]
    \centering
    \includegraphics[width=\linewidth]{../results/figures/cate_by_department.png}
    \caption{Average conditional treatment effects by department.
        Color: direct ATE magnitude; rows sorted by direct effect.
        Departments with fewer than 10 observations are excluded.}
    \label{fig:cate}
\end{figure}

\begin{figure}[t]
    \centering
    \includegraphics[width=\linewidth]{../results/figures/cannibalization_by_dept.png}
    \caption{Substitute spillover (cannibalization) by department.
        Negative values indicate demand loss when substitute neighbors are
        promoted.}
    \label{fig:cannibalization}
\end{figure}

Figures~\ref{fig:cate} and~\ref{fig:cannibalization} decompose effects
by retail department.

\paragraph{Strongest direct responders.}
Salad Bar (ATE $= +0.080$), Garden Center ($+0.074$), Produce ($+0.067$),
and Toys ($+0.057$) show the largest direct responses to promotion.
These are either perishable or discretionary categories where promotional
salience has an outsized effect on purchase incidence.

\paragraph{Weakest or negative direct responders.}
Kiosk/Gas ($-0.026$) and Automotive ($-0.008$) show negative direct
effects.
These likely reflect the fact that fuel and automotive products follow
external price signals (posted pump prices), and store-level promotions
may signal lower quality or confuse customers rather than drive volume.

\paragraph{Cannibalization leaders.}
Travel \& Leisure ($-0.025$), Postal Center ($-0.013$), and Spirits
($-0.009$) show the strongest substitute cannibalization effects.
These are narrowly defined categories where promoted SKUs directly compete
with close alternatives in the same section.

\paragraph{Complement attachment.}
Several departments show positive complement effects in the discount
model (Table~\ref{tab:discount}), consistent with the idea that
deeply discounted staples drive basket expansion.


% ─────────────────────────────────────────────────────────────────────────────
\section{Ablation Study}

To isolate the contribution of each modeling component, we run five
variants of the full model:

\begin{itemize}
    \item \textbf{Full model}: edge-weighted graph + Node2Vec embeddings
    \item \textbf{Binary weights}: replace PPMI/cosine edge weights with
          unweighted adjacency ($w_{ij} \in \{0,1\}$)
    \item \textbf{No embeddings}: omit Node2Vec coordinates from $W$
    \item \textbf{Substitutes only}: zero out $A^{\mathrm{comp}}$
    \item \textbf{Complements only}: zero out $A^{\mathrm{sub}}$
\end{itemize}

Table~\ref{tab:ablation} and Figure~\ref{fig:ablation} summarize the results.

\begin{table}[ht]
    \centering
    \caption{Ablation study results. All variants use the same 6\,M-row
        training sample. ``---'' marks structurally excluded effects.}
    \label{tab:ablation}
    \setlength{\tabcolsep}{4pt}
    \begin{tabular}{lccc}
        \toprule
                         & Direct    & Sub spill. & Comp spill. \\
        \midrule
        Full model       & $+0.0373$ & $-0.0022$  & $+0.0003$   \\
        Binary weights   & $+0.0372$ & $-0.0022$  & $+0.0008$   \\
        No embeddings    & $+0.0431$ & $-0.0048$  & $+0.0024$   \\
        Substitutes only & $+0.0373$ & $-0.0021$  & ---         \\
        Complements only & $+0.0369$ & ---        & $+0.0001$   \\
        \bottomrule
    \end{tabular}
\end{table}

\begin{figure}[t]
    \centering
    \includegraphics[width=\linewidth]{../results/figures/ablation_comparison.png}
    \caption{Ablation study: estimated ATEs with 95\% CIs across five
        model variants. Hatched bars marked ``(excluded)'' indicate
        structurally zeroed-out exposures. The annotation on the Direct
        Effect panel highlights the +16\% inflation when embeddings are removed.}
    \label{fig:ablation}
\end{figure}

\paragraph{Edge weights vs.\ binary.}
The estimates are nearly identical ($\Delta \approx 0.0001$).
This tells us that graph \emph{topology}---which products are neighbors---is
what matters, not the precise calibration of edge weights.
For this dataset, a simpler binary co-purchase graph appears sufficient.

\paragraph{Removing embeddings.}
This is the most consequential ablation.
The direct ATE jumps from $+0.037$ to $+0.043$ ($+16\%$), and the
substitute spillover nearly doubles from $-0.0022$ to $-0.0048$.
Node2Vec coordinates are included in the covariate matrix $W$, where they
act as rich product-level controls.
Without them, product-level confounding that correlates with both promotion
likelihood and baseline demand appears to leak into the treatment effect estimates.
This is analogous to omitted variable bias: categories that are both
more promotable and more elastic look like they have larger treatment effects
unless product-level fixed effects---here, approximated by continuous
embeddings---are controlled.
While embeddings may capture both latent demand structure and residual
confounding, the magnitude and direction of the ablation shift are
consistent with omitted-variable bias at the product level.

\paragraph{Sub-only vs.\ comp-only.}
Direct and substitute effects are stable whether or not the complement
graph is included (and vice versa), confirming that the two spillover
channels are approximately orthogonal in this dataset.
This is expected given that substitute and complement relationships are
built from entirely different co-purchase signals.


% ─────────────────────────────────────────────────────────────────────────────
\section{Synthetic Validation}

To assess whether the DML estimator recovers known effects, we ran a
semi-synthetic simulation study: we kept the real covariate and network
structure but replaced the outcome with a data-generating process with
known direct and spillover coefficients.
We ran 50 independent simulations.
Results are shown in Table~\ref{tab:synthetic}.

\begin{table}[ht]
    \centering
    \caption{Semi-synthetic simulation study: bias, RMSE, and 95\% CI
        coverage over 50 replications.}
    \label{tab:synthetic}
    \begin{tabular}{lcccc}
        \toprule
        Effect         & Mean bias  & RMSE      & Coverage & $n_{\mathrm{sims}}$ \\
        \midrule
        Direct         & $-0.00075$ & $0.00088$ & 0.54     & 50                  \\
        Sub spillover  & $+0.00022$ & $0.00062$ & 0.92     & 50                  \\
        Comp spillover & $-0.00096$ & $0.00119$ & 0.74     & 50                  \\
        \bottomrule
    \end{tabular}
\end{table}

\paragraph{Bias and RMSE.}
All three effects have near-zero mean bias and small RMSE relative to the
true effect sizes, indicating that the estimator recovers the right
parameter values on average.

\paragraph{CI coverage.}
The substitute spillover achieves near-nominal coverage (92\%), consistent
with well-calibrated inference.
The complement spillover falls slightly short (74\%), likely due to the
noisier signal in the complement graph.
The direct effect has notably low coverage (54\%), meaning its 95\%
confidence intervals are too narrow.
This is a known issue with two-way cluster SEs when the primary treatment
has high within-cluster correlation: \texttt{promo\_any} tends to be set
at the product-week level, generating strong intra-product correlation that
shrinks the effective sample size for the direct treatment more than the
asymptotic approximation assumes.
Confidence intervals for the direct effect should therefore be interpreted as
approximate rather than fully calibrated.


% ─────────────────────────────────────────────────────────────────────────────
\section{Discussion}

\paragraph{What graph structure adds.}
The graph matters for causal estimation in two distinct ways.
First, it defines the exposure mappings that make downstream substitute and
complement effects estimable at all.
Second, its embedding representation supplies dense product-level controls
that improve adjustment for confounding.
The ablation results indicate that the second channel is the largest
empirical contributor in this dataset: Node2Vec coordinates act like
approximate product fixed effects, while the precise calibration of edge
weights matters less than the neighborhood topology itself.

\paragraph{Why complement spillovers are weak.}
The complement attachment effect is consistently positive but never
statistically significant in the binary promotion model, though it does
reach significance in the continuous discount model.
This pattern suggests a substantive distinction between extensive-margin and
intensive-margin promotion measures: binary promotions (display/mailer flags)
are noisy indicators of price salience, while a deeper discount more directly
captures the price shock that triggers basket expansion.
Another possibility is that one-hop complement aggregation is too coarse
to capture the full attachment signal, which may propagate through longer
purchase paths.

\paragraph{Limitations.}
This study has several limitations worth noting:

\begin{itemize}
    \item \textbf{One-hop exposure only.} The exposure mapping in
          Eq.~\eqref{eq:exposure} aggregates only immediate graph neighbors.
          Multi-hop propagation is not captured.

    \item \textbf{Static product graphs.} Both adjacency matrices are built
          from the full transaction history, ignoring seasonal or assortment
          shifts in product relationships.

    \item \textbf{Observational data.} Despite the DML design, unmeasured
          store-week confounders (local events, weather, competitor promotions)
          could bias spillover estimates if they affect multiple products
          simultaneously.
          In particular, coordinated store-week promotions affecting multiple
          related products could induce correlated treatment assignment and
          outcomes not fully captured by observed controls.

    \item \textbf{Direct effect CI undercoverage.} As noted in
          Section~VIII, the 95\% CIs for the direct effect cover the true value
          only 54\% of the time in simulation, suggesting that significance
          claims for the direct ATE should be interpreted conservatively.

    \item \textbf{Treatments estimated separately.} Binary \texttt{promo\_any}
          and continuous \texttt{discount\_rate} are estimated in separate models.
          Joint estimation of both treatment dimensions is left for future work.
\end{itemize}


% ─────────────────────────────────────────────────────────────────────────────
\section{Conclusion}

We presented a graph-aware Double Machine Learning framework for estimating
causal spillovers in retail pricing.
The key design choices are: (1) representing product relationships as a
weighted graph built from transaction co-occurrence and embedding similarity;
(2) computing one-hop weighted exposure variables to summarize neighborhood
treatment intensity; and (3) embedding product coordinates in the covariate
matrix as approximate product fixed effects.

On the Dunnhumby retail panel, we estimate a positive direct promotion effect
of $+0.037$ on the $\log(1 + \text{units sold})$ scale, a statistically
significant negative substitute spillover, and a directional complement
attachment effect that becomes significant under continuous discount depth.
Importantly, the ablation study demonstrates that product embeddings are
critical: without them, the direct effect rises by 16\% and the substitute
spillover estimate more than doubles.
More broadly, the results show that incorporating product-relationship
structure improves estimation of downstream treatment effects in this retail
setting by making spillovers explicit and by providing richer controls for
product-level confounding.

These results suggest that retailers should account for cross-product demand
spillovers when evaluating promotional strategies, and that embedding-based
product representations are a practical and effective way to control for
product-level confounding in large-scale observational datasets.


% ─────────────────────────────────────────────────────────────────────────────
\bibliographystyle{IEEEtran}
\bibliography{refs}

\end{document}
